% §3 Security Framework for Covenant Vault Designs
% Structure: per-class subsections (Z1–Z6), each deriving
% susceptibility and immunity from the axes of §2.3. Every proof
% preserved verbatim from the prior axis-structured draft.

\section{Security Framework for Covenant Vault Designs}
\label{sec:impossibility}

This section develops the security framework used throughout the
paper. Section~\ref{sec:imposs:framework} defines the defender-loss
metric that unifies theft, griefing, and dust-lockup outcomes.
Sections~\ref{sec:imposs:class_fee}--\ref{sec:imposs:class_mode}
develop six attack classes \classfee{}--\classmode{}; each class is a
mechanical consequence of a specific axis-value from
Table~\ref{tab:design_space}, with susceptibility and immunity
derived side-by-side rather than asserted.
Section~\ref{sec:imposs:design_space} identifies a structurally
dominant corner that no Bitcoin proposal occupies.
Section~\ref{sec:imposs:sensitivity} applies the framework
threat-by-threat, characterising how defender loss scales with fee
rate and recovery strategy.

% ================================================================
\subsection{Defender Loss and Adversary Advantage}
\label{sec:imposs:framework}

We parameterise each scenario by a threat model
$\mathrm{TM} \in \mathcal{T}$, design $D \in \mathcal{D}$, vault
value $V > 0$, fee rate $r \geq 0$, and defender strategy
$b \in \{0, 1\}$ ($b = 1$: defender batches recoveries; $b = 0$:
single-input). An adversary strategy $\mathcal{A}$ is a sequence of
on-chain transactions in pursuit of $\mathrm{TM}$.

\begin{definition}[Costs and outcomes]
\label{def:strategy}
For strategy $\mathcal{A}$ at fee rate $r$: $g(\mathcal{A}) \in [0,V]$
are extracted funds, $u(\mathcal{A}) \in [0,V]$ are dust-locked funds,
$c_A(\mathcal{A}, r)$ the adversary's fees, and $c_D(\mathcal{A}, r, b)$
the defender's fees under optimal response.
\end{definition}

\begin{definition}[Defender loss and adversary advantage]
\label{def:loss-adv}
Fix $(\mathrm{TM}, D, V, r, b)$ and attacker budget $B_A = V$
(rational-adversary assumption). The \emph{worst-case defender loss}
and \emph{adversary advantage} are
\[
  \mathcal{L} = \!\!\sup_{c_A(\mathcal{A}, r) \leq B_A}\!
    [\,g(\mathcal{A}) + u(\mathcal{A}) + c_D(\mathcal{A}, r, b)\,],
  \qquad
  \mathrm{Adv} = \sup_{\mathcal{A}} [\,g(\mathcal{A}) - c_A(\mathcal{A}, r)\,].
\]
$D_1 \prec_{\mathrm{TM}, V, r, b} D_2$ iff
$\mathcal{L}(\mathrm{TM}, D_1, V, r, b) < \mathcal{L}(\mathrm{TM}, D_2, V, r, b)$.
$\mathcal{L}$ unifies theft ($g>0$), griefing ($c_D > 0$), and
dust-lockup ($u > 0$); $\mathrm{Adv}$ is the attacker's rationality
signal and is $b$-invariant.
\end{definition}

Every covenant vault design is characterised by four independent
formal choices: a \emph{fee model}~$f$ (axis \axfee{}), an
\emph{amount-flexibility policy}~$a$ (axis \axrevault{}), a
\emph{recovery gating}~$g$ (axis \axauth{}), and a \emph{recovery
binding}~$b$ (jointly axes \axwdbind{} and \axrecbind{}). We write
$D = (f, a, g, b)$. Two spec-level dependencies link the axes: a fee
model that commits to a single output
(\sighash{SINGLE}|\texttt{ANYONECANPAY}) forces $a_{\text{atomic}}$;
conversely, $a_{\text{partial}}$ requires per-output introspection and
therefore $f \in \{f_{\text{wallet}}, f_{\text{inval}}\}$.
Table~\ref{tab:design_space} (Section~\ref{sec:bg:axes}) lists each
design's position on every axis; the six subsections that follow
derive the attack class enabled at each axis-value.

% ================================================================
\subsection{Class \classfee{}: Fee-Channel Pinning}
\label{sec:imposs:class_fee}

\paragraph{Class definition.} \classfee{} is the class enabled by
axis \axfee{} = \texttt{out-of-band-anchor}: fees attach via a CPFP
anchor output, whose own unconfirmed descendants block
defender-initiated replacement under standard-policy mempool rules.

\paragraph{The four fee models.} Four spec-level choices appear in
the proposals we study. \emph{CPFP anchor} ($f_{\text{anchor}}$,
\opctv{}): the unvault carries a small externally-spendable output
(${\approx}3{,}300$~sats) that lets any party bump the fee via
Child-Pays-For-Parent. \emph{Fee-wallet input} ($f_{\text{wallet}}$,
\opvault{}): value preservation forbids consuming vault value for
fees, so every trigger and recovery transaction attaches a separate
P2TR input from a fee wallet (adding 80--90\vB{}). \emph{In-value
deduction} ($f_{\text{inval}}$, \opccv{}): per-output introspection
allows fees to be deducted from the vault value itself. \emph{Sighash
ACP} ($f_{\text{acp}}$, \catcsfs{}): \sighash{SINGLE}|\texttt{ANYONECANPAY}
commits the signature to a single output and lets external
fee-paying inputs attach without breaking the covenant.

\paragraph{Susceptibility (implication D1).} \opctv{} is susceptible.
Under Bitcoin Core's default 25-descendant mempool relay policy, a
hot-key adversary broadcasts the unvault and attaches $25$ dust-rate
transactions spending the anchor. BIP-125 Rule 3 requires the
watchtower's CPFP replacement to pay for evicting every descendant
at a higher fee rate and higher absolute fee; against the attacker's
dust chain this is economically infeasible, blocking recovery and
forcing $\mathcal{L}(\mathrm{TM1}) = V$.

\paragraph{Immunity derivations.} \opccv{}, \opvault{}, and \catcsfs{}
have $\axfee{} \neq \texttt{out-of-band-anchor}$ and therefore carry
no pinnable anchor handle: dynamic fee inclusion keeps fee payment
inside the transaction the attacker would need to replace. The
attacker cannot pin a transaction whose existence they are trying to
delay. $\mathcal{L}(\mathrm{TM1}) = 0$ for all three.

\paragraph{Rationality (Appendix B.2).} \emph{Attacker capability:}
hot key (Kerckhoffs, §\ref{sec:bg:lifecycle}); capital for the
anchor-descendant chain ($\approx 14{,}300$~sats at $\rho = 3\satvB$).
\emph{Defender model:} online watchtower, Core-default mempool
policy. \emph{Preconditions:} $\mathcal{L}$ and $\mathrm{Adv}$ are
$r$-invariant. \emph{Counter-factual:} TRUC/BIP-431 collapses the
descendant limit to $1$, eliminating the pinning
surface~\cite{TRUC}; cluster mempool~\cite{ClusterMempool}
qualitatively changes the replacement calculus but is a scope-limit
rather than a within-scope mitigation. Full block:
Appendix~\ref{app:rationality}.

\paragraph{Empirical.} Section~\ref{sec:empirical:classes} reports
$C_{\text{pin}}$ and attack-composition with TM3.

% ================================================================
\subsection{Class \classgrief{}: Permissionless Griefing}
\label{sec:imposs:class_grief}

\paragraph{Class definition.} \classgrief{} is the class enabled by
axis \axauth{} = \texttt{keyless}: emergency recovery is invocable
without a signature, so any party --- including a pure griefer ---
can cycle the owner through re-deposit and re-trigger.

\paragraph{The two gating models.} Emergency recovery can be
\emph{keyless} ($g_{\text{keyless}}$, \opccv{}), invocable by any
party, or \emph{key-gated} ($g_{\text{key}}$,
\opctv{}\,/\,\opvault{}\,/\,\catcsfs{}), requiring a signature under a
specific cold or authorisation key. We formalise this choice via an
authorisation function
$\mathcal{A}\colon 2^{\mathcal{K}} \to \{0,1\}$ over the key set
$\mathcal{K}$: recovery is permissionless iff
$\mathcal{A}(\emptyset) = 1$.

\begin{proposition}[Griefing--safety incompatibility]
\label{thm:griefing_safety}
No covenant vault simultaneously achieves (i) permissionless
recovery ($\mathcal{A}(\emptyset) = 1$) and (ii) griefing
resistance (no unauthorised party can invoke recovery to force
$c_D > 0$).
\end{proposition}

\begin{proof}
Case analysis on $\mathcal{A}(\emptyset)$. If
$\mathcal{A}(\emptyset) = 1$, any third party meets the
authorisation predicate and can force recovery, requiring the
owner to re-deposit and re-trigger at cost $c_D > 0$; (ii) is
violated. If $\mathcal{A}(\emptyset) = 0$, unauthorised parties
cannot invoke recovery so (ii) holds, but (i) is violated by
definition. The cases exhaust
$\{0, 1\}$. \qed
\end{proof}

\paragraph{Susceptibility (implication D3).} \opccv{} is susceptible
and is the only design with \axauth{} = \texttt{keyless}. The
attacker pool is the entire population of Bitcoin nodes that can
observe the mempool; any such node can construct and broadcast a
valid recovery transaction consuming the unvaulting UTXO. The
Kerckhoffs assumption makes the recovery address public (cold
addresses are routinely disclosed for auditability or derived by the
same open-source library), so taproot leaf hiding does not provide
security here. $\mathcal{L}(\mathrm{TM2}) > 0$ for any fee rate.

\paragraph{Immunity derivations.} \opvault{} has \axauth{} =
\texttt{recoveryauth-key}: recovery demands a Schnorr signature
under a dedicated key, reducing the attacker pool to the key holder
and their compromises. \opctv{} has \axauth{} = \texttt{hot-key}
(recovery is the cold sweep, signed by the hot key with
\opctv{}-committed destination); the griefing surface collapses once
the hot key is safe. \catcsfs{} has \axauth{} = \texttt{cold-key}:
griefing requires cold-key compromise, which the attacker would
prefer to convert into class \classcold{} (Section~\ref{sec:imposs:class_cold}).

\paragraph{Rationality (Appendix B.2).} \emph{Attacker capability:}
zero keys; observation of mempool plus the public recovery address.
\emph{Defender model:} any; the attack is liveness-only and ends
only when the defender abandons or rotates the vault.
\emph{Preconditions:} $\mathcal{L}$ scales linearly in $r$;
$\gamma = \text{vsize}_{\text{atk}}/\text{vsize}_{\text{def}}$ is
$r$-invariant. \emph{Counter-factual:} upgrading \axauth{} to
\texttt{recoveryauth-key} (add a required signature on the recovery
leaf) eliminates the class; this is BIP-345's design choice.

\paragraph{Empirical.} Section~\ref{sec:empirical:classes} reports
$\gamma$ per design.

% ================================================================
\subsection{Class \classscale{}: Per-UTXO Recovery-Cost Scaling}
\label{sec:imposs:class_scale}

\paragraph{Class definition.} \classscale{} is the class enabled by
axis \axrevault{} = \texttt{yes}: partial withdrawal grows the
vault's UTXO count over time. Per-UTXO recovery cost scales linearly
in count and in fee rate; under congestion, some UTXOs fall below
the economic abandonment threshold. An attacker with the trigger
key can accelerate the distribution; ordinary operation can produce
it without one.

\paragraph{The two amount policies.} \emph{Atomic}
($a_{\text{atomic}}$, \opctv{}, \catcsfs{}): the full vault value
must exit the covenant in one transaction. \emph{Partial with
revault} ($a_{\text{partial}}$, \opccv{}, \opvault{}): per-output
introspection lets the spender produce a withdrawal portion and a
re-locked vault UTXO atomically, consuming one vault UTXO and
creating two outputs that between them preserve the original value.

\paragraph{Susceptibility (implication D2).} \opccv{} and \opvault{}
are susceptible. Repeated partial unvaults fragment the vault into
progressively smaller UTXOs. Each UTXO the defender must recover
costs $\text{vsize}_{\text{rec}} \cdot r$ in fees; when this exceeds
the UTXO's value, recovery is economically irrational and the UTXO
strands as dust. Above a fee-rate threshold derived in
Section~\ref{sec:imposs:sensitivity}, the defender loss
$\mathcal{L}(\mathrm{TM4})$ is strictly positive.

\paragraph{Immunity derivations.} \opctv{} and \catcsfs{} have
$a_{\text{atomic}}$ and cannot express partial withdrawal: the
template hash of \opctv{} and the sighash-preimage binding of
\catcsfs{} both commit to a fixed output set with no provision for
splitting. $\mathcal{L}(\mathrm{TM4}) = 0$ at every fee rate.

\paragraph{Reframing note.} Prior framings of this class as
``watchtower exhaustion'' cast it as a multi-round attacker-versus-
watchtower game. That framing does not survive a defender with full
recovery authority, who can full-sweep on first alert. The reframe
above identifies the class as a structural cost tail of partial
withdrawal --- attacker-accelerable but not attacker-required --- and
locates the defender's degree of freedom in recovery batching
(\ref{sec:empirical:batching}).

\paragraph{Rationality (Appendix B.2).} \emph{Attacker capability:}
trigger (hot) key (Kerckhoffs); capital for $N$ partial-withdrawal
triggers. \emph{Defender model:} watchtower with recovery authority;
batching flag $b$ governs per-UTXO cost scaling; minimal-disruption
policy (reactive per-UTXO recovery) maximises attacker leverage.
\emph{Preconditions:} rate-dependent; threshold
$r^\dagger_{\mathrm{TM4}}$ defined in
Section~\ref{sec:imposs:sensitivity}. \emph{Protocol assumptions:}
Bitcoin dust thresholds ($546$~sats P2WPKH, $330$~sats P2TR); RBF
rule 3 does not permit sibling-eviction, so the attacker's
partial-withdrawal chain is single-chain.
\emph{Counter-factual:} defender batching raises $r^\dagger$ by
$\mathord{\approx}\!1.8$--$2.7\times$
(Section~\ref{sec:empirical:batching}), absorbing the tail under
moderate congestion; full-sweep-on-first-alert policy collapses the
class to one round at the cost of operational disruption.

\paragraph{Empirical.} Section~\ref{sec:empirical:batching} reports
$(\alpha, \beta, r^\dagger)$ per design and batching flag, including
the \opccv{}/\opvault{} ordering flip that is the paper's novel
finding.

% ================================================================
\subsection{Class \classhot{}: Hot-Key Theft via Destination Substitution}
\label{sec:imposs:class_hot}

\paragraph{Class definition.} \classhot{} is the class enabled by
axis \axwdbind{} $\neq$ \texttt{signature-dual-bind}: the hot-key
holder, upon compromise, can direct the withdrawn funds to an
attacker-chosen destination, racing the watchtower during the CSV
delay. Signature-dual-binding on the withdrawal path uniquely
forecloses the redirection primitive at the cryptographic layer.

\paragraph{Susceptibility.} \opctv{} (\axwdbind{} =
\texttt{template-hash-digest}), \opccv{} (\axwdbind{} =
\texttt{per-output-amount+script}), and \opvault{} (\axwdbind{} =
\texttt{dedicated-opcode}) are all susceptible. In each case the
hot-key holder's signature authorises an unvault; the destination is
constrained by the \emph{script path spent} but the hot key can
select among alternative leaves or construct an unvault whose
destination the defender must recover from under the CSV race. The
outcome depends on output binding at unvault:
$\mathcal{L} = V \cdot \Pr_D[\text{attacker wins}] +
c_D^{\text{recover}} \cdot \Pr_D[\text{defender wins}]$.
$\Pr_D$ is dominated by $r$-independent factors (CSV length, network
topology), except that $\Pr_{\mathrm{OP\_CTV}}$ rises with TM1
applicability: class \classfee{} composes with \classhot{} under
\opctv{} to give $\mathcal{L} = V$ when fee pinning succeeds.

\paragraph{Immunity derivation: \catcsfs{}.} \catcsfs{} is immune by
\axwdbind{} = \texttt{signature-dual-bind}. Let $\sigma$ be the
Schnorr signature supplied in the witness, $m_{\text{preimage}}$ the
stack-assembled BIP-341 sighash preimage, and $m_{\text{consensus}}$
the sighash computed by consensus over the actual transaction.
\opcsfs{} checks
$\mathrm{Verify}(\mathrm{hot\_pk},\,m_{\text{preimage}},\,\sigma) = 1$;
\texttt{OP\_CHECKSIG} checks
$\mathrm{Verify}(\mathrm{hot\_pk},\,m_{\text{consensus}},\,\sigma) = 1$.
Both hold only if $m_{\text{preimage}} = m_{\text{consensus}}$: a
valid signature over two distinct messages under the same key would
violate Schnorr's EUF-CMA security~\cite{NRS20}, reducing to
discrete-log. The \texttt{hashOutputs} field is a component of both
preimages, so the output set is forced. Attempting to redirect yields
a rejected transaction, not a theft; the outcome degrades to
griefing-only (TM9). $\mathcal{L} = c_D$.

\paragraph{Rationality (Appendix B.2).} \emph{Attacker capability:}
hot key; mempool access; optional fee-pinning capital (if \classfee{}
composes). \emph{Defender model:} online watchtower able to race
during CSV; race bias toward attacker under TM1 composition.
\emph{Preconditions:} CSV length, network topology; $r$-independent
except via TM1.
\emph{Counter-factual:} \axwdbind{} = \texttt{signature-dual-bind}
(as in \catcsfs{}) eliminates the class by cryptographic construction;
any covenant able to express dual-verification can port the
primitive.

\paragraph{Empirical.} Section~\ref{sec:empirical:classes} reports
$\mathcal{L}$ per design and the composition with \classfee{}.

% ================================================================
\subsection{Class \classcold{}: Cold-Key Theft via Unconstrained Recovery}
\label{sec:imposs:class_cold}

\paragraph{Class definition.} \classcold{} is the class enabled by
axis \axrecbind{} = \texttt{plain-signature}: the recovery path is a
bare \texttt{cold\_pk OP\_CHECKSIG} with no script-level
output constraint. Cold-key compromise produces unrestricted fund
redirection rather than forced recovery to the pre-committed
destination.

\paragraph{The two recovery-binding models.} \emph{Output-bound}
($b_{\text{bound}}$, \opctv{}, \opvault{}, \opccv{}, Simplicity): the
recovery output scriptPubKey is fixed at vault creation (via a
\opctv{} template, a \opvault{}/\opccv{}-committed destination, or a
Simplicity output-hash jet). \emph{Unbound}
($b_{\text{unbound}}$, \catcsfs{} recovery leaf): the recovery leaf
is a bare \texttt{cold\_pk OP\_CHECKSIG} with \sighash{DEFAULT},
placing no constraint on outputs.

\paragraph{Susceptibility (implication D4).} \catcsfs{} is
susceptible. A compromised cold key produces a valid signature over
any sighash, which \texttt{OP\_CHECKSIG} accepts; there is no
script-level check on \texttt{hashOutputs}.
$\mathcal{L}(\mathrm{TM11}) = V$; the attacker chooses the
destination freely. $r$-invariant.

\paragraph{Immunity derivations.} \opctv{} and \opvault{} have
\axrecbind{} \,$\in$\, \{\texttt{template-hash-digest},
\texttt{dedicated-opcode}\}: the recovery script commits to a
specific recovery destination, so a compromised cold key can force
recovery but cannot redirect funds.
$\mathcal{L}(\mathrm{TM11}) = 0$. \opccv{} is excluded by
\axauth{} = \texttt{keyless} (no cold key to compromise); its
\axrecbind{} = \texttt{per-output-amount+script} would likewise
prevent redirection.

\paragraph{Note (axis \axwdbind{} vs.\ \axrecbind{} asymmetry).}
\catcsfs{}'s strongest axis position is \axwdbind{} =
\texttt{signature-dual-bind} (uniquely immune to \classhot{}) paired
with its weakest axis position \axrecbind{} = \texttt{plain-signature}
(uniquely susceptible to \classcold{}). Splitting $b$ into the two
sub-axes \axwdbind{} and \axrecbind{} (§\ref{sec:bg:axes}) turns
this asymmetry from a remark into a mechanical reading of the
position table. \cite{Poe21} sketches a \catcsfs{}-based bound
recovery leaf that would move the design to \axrecbind{} =
\texttt{signature-dual-bind}, eliminating \classcold{}.

\paragraph{Rationality (Appendix B.2).} \emph{Attacker capability:}
cold key (Kerckhoffs --- the attack assumes compromise, not
inference). \emph{Defender model:} any; the attack is instant with no
CSV window. \emph{Counter-factual:} \axrecbind{} =
\texttt{signature-dual-bind} on the recovery leaf eliminates the
class; alternatively, \axauth{} = \texttt{keyless} prevents the
cold-key compromise surface but is Pareto-dominated by Proposition~1.

\paragraph{Empirical.} Section~\ref{sec:empirical:classes}.

% ================================================================
\subsection{Class \classmode{}: Mode/Parameter Bypass}
\label{sec:imposs:class_mode}

\paragraph{Class definition.} \classmode{} is the class enabled by
axis \axopsurf{} = \texttt{multi-mode-with-OP\_SUCCESS}: a
polymorphic covenant opcode takes a stack-provided mode selector,
and the taproot \texttt{OP\_SUCCESSx} upgrade mechanism promotes
undefined mode values to unconditional script success. A developer
or construction bug that selects an undefined mode silently disables
the covenant check.

\paragraph{Susceptibility.} \opccv{} is susceptible. BIP-443 defines
four valid modes $\{-1, 0, 1, 2\}$; any other value falls through to
\texttt{OP\_SUCCESSx}, which under BIP-341/342 taproot rules
terminates script execution with success regardless of stack state.
Five undefined modes $\{3, 4, 7, 128, 255\}$ were measured on
regtest; each produces full covenant bypass at $110$\vB{} per spend.
$\mathcal{L}(\mathrm{TM8}) = V$.

\paragraph{Immunity derivations.} \opctv{}, \opvault{}, and \catcsfs{}
have \axopsurf{} = \texttt{single-mode}: each exposes a single
semantic opcode with no mode selector, so the \texttt{OP\_SUCCESSx}
fallthrough surface does not arise.

\paragraph{Framing.} The class is \emph{specified behaviour} under
BIP-443's forward-compatibility design (itself inherited from
BIP-341/342), not a newly-discovered vulnerability. The contribution
here is the systematic mode sweep and the class-level framing that
ties the surface to axis \axopsurf{}, not the discovery of
\texttt{OP\_SUCCESSx} semantics. In practice, the pymatt reference
library mitigates the risk by providing named constants
(\texttt{CCV\_FLAG\_*}); the residual risk is in independent
implementations that hand-write the mode byte.

\paragraph{Rationality (Appendix B.2).} \emph{Attacker capability:}
any observer of a vault whose script contains an undefined mode
byte; no keys required. \emph{Defender model:} developer discipline
around mode-byte selection; library-provided named constants.
\emph{Counter-factual:} \axopsurf{} = \texttt{single-mode} (one
opcode per primitive, as in \opctv{}, \opvault{}, \catcsfs{})
eliminates the class. A bitmask-gated version of \opccv{} (reject
undefined modes at script-parse time) would also eliminate the
class without requiring a soft-fork revision, but would forfeit the
forward-compatibility benefit.

\paragraph{Empirical.} Appendix~\ref{app:threat_matrix} documents the
five-mode sweep.

% ================================================================
\subsection{The Dominant Corner and Current Proposals}
\label{sec:imposs:design_space}

Combining the axes, each design occupies one cell of the
lattice. Restricting to the proof-relevant formal sublattice
$(f, a, g, b)$, define the \emph{dominant corner}
$D^\star = (f_{\neq\text{anchor}}, a_{\text{atomic}},
g_{\text{key}}, b_{\text{bound}})$. By implications D1--D4, any
design at $D^\star$ has $\mathcal{L} = 0$ simultaneously against
\{TM1, TM2, TM4, TM11\} --- equivalently, against classes
\{\classfee{}, \classgrief{}, \classscale{}, \classcold{}\}.
Table~\ref{tab:pareto-summary} locates each proposed BIP in the
sublattice with Hamming distance to $D^\star$; the full axis
positions are Table~\ref{tab:design_space} in
Section~\ref{sec:bg:axes}.

\begin{table}[t]
\centering
\caption{The four Bitcoin Script extensions in the $(f, a, g, b)$
sublattice of Table~\ref{tab:design_space}. Distance measures Hamming
separation from the dominant corner $D^\star$ on the proof-relevant
sublattice; the full A1--A7 position of each design is available in
Section~\ref{sec:bg:axes}.}
\label{tab:pareto-summary}
\small
\begin{tabular}{@{}lccccc@{}}
\toprule
Design     & $f$                      & $a$                  & $g$                     & $b$                    & dist.\ \\
\midrule
\opctv{}        & $f_{\text{anchor}}$      & $a_{\text{atomic}}$  & $g_{\text{key}}$        & $b_{\text{bound}}$     & 1 \\
\opccv{}        & $f_{\text{inval}}$       & $a_{\text{partial}}$ & $g_{\text{keyless}}$    & $b_{\text{bound}}$     & 2 \\
\opvault{} & $f_{\text{wallet}}$      & $a_{\text{partial}}$ & $g_{\text{key}}$        & $b_{\text{bound}}$     & 1 \\
\catcsfs{}   & $f_{\text{acp}}$         & $a_{\text{atomic}}$  & $g_{\text{key}}$        & $b_{\text{unbound}}$   & 1 \\
$D^\star$  & $\neq f_{\text{anchor}}$ & $a_{\text{atomic}}$  & $g_{\text{key}}$        & $b_{\text{bound}}$     & --- \\
\bottomrule
\end{tabular}
\end{table}

No proposed Bitcoin BIP occupies $D^\star$. \opctv{} differs
in~$f$ and inherits \classfee{} by D1; \opccv{} differs in~$a$
and~$g$ and inherits \classscale{} and \classgrief{} by D2 and D3;
\opvault{} differs in~$a$ and inherits \classscale{} by D2;
\catcsfs{} differs in~$b$ and inherits \classcold{} by D4. A
hypothetical specification at $D^\star$ would have $\mathcal{L} = 0$
against all four of \{\classfee{}, \classgrief{}, \classscale{},
\classcold{}\} simultaneously; the corner is currently open on
Bitcoin. Two present proposals are one axis away: \opctv{}
(by replacing its anchor-based fee model, which TRUC
adoption~\cite{TRUC} approximates) and \catcsfs{} (by replacing its
bare cold-key recovery leaf with an output-bound path, sketched
in~\cite{Poe21}). Simplicity on Elements (Appendix~D) occupies the
dominant corner --- demonstrating the corner is \emph{occupiable},
just not by any current Bitcoin proposal.
\label{thm:fee_inversion}\label{cor:open_corner}

\paragraph{Localisation for future proposals.}
The framework yields a lookup procedure: read a proposal's
specification, extract its axis tuple
$(\axauth, \axrevault, \axfee, \axwdbind, \axrecbind, \axopsurf, \axintro)$,
and consult the class derivations
(\ref{sec:imposs:class_fee}--\ref{sec:imposs:class_mode}) to read
off its vulnerability profile. No further measurement is needed to
characterise which of \{\classfee{}, \classgrief{}, \classscale{},
\classhot{}, \classcold{}, \classmode{}\} a newly-proposed BIP
admits by construction. What the framework does \emph{not}
determine---numerical thresholds like $r^\dagger_{\mathrm{TM4}}(V,
D, b)$---depends on user parameters ($V$, watchtower block budget)
and measured transaction sizes, and is the subject of
Section~\ref{sec:imposs:sensitivity} and the empirical measurements
of Section~\ref{sec:empirical}.

% ================================================================
\subsection{Per-Threat Loss Characterisation}
\label{sec:imposs:sensitivity}

We now apply the framework threat by threat, deriving closed-form
expressions for $\mathcal{L}$ as a function of fee rate and
defender strategy. Each TM is the scope-anchor of the class noted in
parentheses; see Appendix~\ref{app:rationality} for the full
seven-field rationality blocks.

\paragraph{TM1 (fee-channel pinning, class \classfee{})~\cite{Har24}.}\label{sec:sens:tm1}
Applies to \opctv{} alone by D1.
$\mathcal{L}(\mathrm{TM1}, \mathrm{OP\_CTV}, V, r) = V$;
$\mathrm{Adv} = V - C_{\text{pin}}$ with
$C_{\text{pin}} = 25(\rho+1)\,\text{vsize}_{\text{desc}} + \alpha_{\text{anchor}}
\approx 14{,}300$~sats ($\rho = 3\satvB$, anchor $= 3{,}300$~sats).
Both $\mathcal{L}$ and $\mathrm{Adv}$ are $r$-invariant;
$\mathcal{L} = 0$ for $D \neq \mathrm{OP\_CTV}$.

\paragraph{TM2 (recovery griefing, class \classgrief{}).}
Applies to all designs, with attacker pools set by D3. Over $k$
rounds,
\begin{equation} \label{eq:gamma}
  \mathcal{L}(\mathrm{TM2}, D, V, r) = k \cdot \text{vsize}_{\text{def}}(D) \cdot r,
  \qquad
  \gamma(D) = \frac{\text{vsize}_{\text{atk}}(D)}{\text{vsize}_{\text{def}}(D)}.
\end{equation}
$\mathcal{L}$ scales linearly in $r$; $\gamma$ is $r$-invariant.
Ordering by $\gamma$: \opccv{} ($0.27$, keyless, severe) $<$ \catcsfs{}
($0.36$) $<$ \opvault{} ($0.55$) $<$ \opctv{} ($0.71$).

\paragraph{TM3 (trigger-key theft, class \classhot{}).}
Applies to all designs. The hot-key adversary unvaults and races
the watchtower during the CSV delay; outcome depends on output
binding at unvault:
$\mathcal{L} = V \cdot \Pr_D[\text{attacker wins}] +
c_D^{\text{recover}} \cdot \Pr_D[\text{defender wins}]$. $\Pr_D$ is
dominated by $r$-independent factors (CSV length, network
topology), except that $\Pr_{\mathrm{OP\_CTV}}$ rises with TM1
applicability. Ordering: \catcsfs{} ($\mathcal{L} = c_D$,
dual-verification binding on the hot leaf) $<$ \opvault{} $\approx$
\opccv{} ($\mathcal{L} \in [c_D, V]$) $<$ \opctv{} ($\mathcal{L} = V$ when
TM1 succeeds).

\paragraph{TM4 (per-UTXO recovery-cost scaling, class \classscale{}).}\label{sec:sens:tm4}
Applies to \opccv{} and \opvault{} by D2. Let $\text{vsize}_{\text{rec}}(D)$
be the single-input recovery vsize. Under batched recovery, BIP-345
§``Batching'' and BIP-443 default aggregation permit the defender
to recover $N$ UTXOs in one transaction of vsize
$\alpha(D) + \beta(D) \cdot N$. The effective per-UTXO recovery
cost is
\begin{equation}
  \text{vsize}_{\text{rec}}(D, b, N) =
  \begin{cases}
    \text{vsize}_{\text{rec}}(D)                             & \text{if } b = 0, \\
    \alpha(D)/N + \beta(D)                                   & \text{if } b = 1.
  \end{cases}
\end{equation}
A split becomes dust when its value falls below this effective
cost; the number of splits needed to drive every UTXO below dust
is $N_{\text{dust}}(D, r, b) = \lfloor V /
(\text{vsize}_{\text{rec}}(D, b, N_{\text{dust}}) \cdot r) \rfloor$.
Let $S(D)$ denote splits per block. The single-block-feasibility
threshold
\begin{equation} \label{eq:r-dagger}
  r^\dagger_{\mathrm{TM4}}(V, D, b)
  \;=\; \frac{V}{\text{vsize}_{\text{rec}}(D, b, \cdot) \cdot S(D)}
\end{equation}
is the fee rate at which $N_{\text{dust}}/S = 1$; lower
$r^\dagger$ means the structural cost tail saturates at lower fees,
hence less safe. Batching strictly raises $r^\dagger$ whenever
$\beta(D) < \text{vsize}_{\text{rec}}(D)$, always true in practice.
Measured $(\alpha, \beta)$ values, the resulting thresholds, and
the \opccv{}/\opvault{} ordering (which flips between $b = 0$ and
$b = 1$) appear in Section~\ref{sec:empirical:batching}. The
multi-round attacker narrative of prior framings is one
threat-model instantiation among several; the class-level statement
is that the cost tail exists as a consequence of $a_{\text{partial}}$
independently of adversary behaviour.

\paragraph{TM11 (cold-key compromise, class \classcold{}).}\label{sec:sens:tm11}
Applies to key-gated designs by D4.
$\mathcal{L}(\mathrm{TM11}) = V$ under $b_{\text{unbound}}$
(\catcsfs{}: bare \texttt{cold\_pk OP\_CHECKSIG} with
\sighash{DEFAULT});
$\mathcal{L}(\mathrm{TM11}) = 0$ under $b_{\text{bound}}$ (\opctv{},
\opvault{}). \opccv{} is excluded by $g_{\text{keyless}}$ (no cold key
to compromise). $r$-invariant in both cases.

\paragraph{Protocol-level vs.\ user-parameter-level.}
The \emph{existence} of the threshold $r^\dagger_{\mathrm{TM4}}$
is a protocol-level consequence of D2. Its numerical value depends
on $V$ (user-chosen vault value), $N_{\text{budget}}$ (defender
block budget), and $\text{vsize}_{\text{rec}}$ fixed by BIP
specification; implementation variation is below 5\vB{} in our
measurements. The implications D1--D4 thus establish
protocol-level structural properties; reported numerical values
scale linearly with user parameters.

\begin{table}[t]
\centering
\caption{Per-class worst-case defender loss $\mathcal{L}$ across
the four Bitcoin covenant proposals. Values are for
$V = 0.5$~BTC. Simplicity is treated separately (Appendix~D).
Every zero-loss cell is annotated with the structural reason,
keyed to the footnote below.
\textbf{Legend:} $V$ = full theft; $0$ = design immune by the
structural property noted;
$\gamma = \text{vsize}_{\text{atk}} / \text{vsize}_{\text{def}}$;
$[c_D, V]$ = race-dependent bounded interval;
$r^\dagger$ = single-block-feasibility fee rate (\satvB{}).
Fee-rate dependence: inv = invariant, $\propto r$ = linear,
thr = threshold, race = race-dependent.}
\label{tab:pareto-matrix}
\renewcommand{\arraystretch}{1.1}
\begin{tabular}{@{}lccccc@{}}
\toprule
Threat (class)                                     & $r$-dep.    & \opctv{}                    & \opccv{}                  & \opvault{}       & \catcsfs{}               \\
\midrule
TM1 pinning (\classfee{})                          & inv         & $V$                    & $0^{\mathrm{a}}$     & $0^{\mathrm{b}}$ & $0^{\mathrm{c}}$       \\
TM2 griefing (\classgrief{})                       & $\propto r$ & $\gamma{=}0.71$        & $\gamma{=}0.27$      & $\gamma{=}0.55$  & $\gamma{=}0.36^{*}$    \\
TM3 trig-theft (\classhot{})                       & race        & $[c_D, V]^{\dagger}$   & $[c_D, V]$           & $[c_D, V]$       & $c_D$                  \\
TM4 cost scaling (\classscale{})                   & thr         & $0^{\mathrm{d}}$       & $r^\dagger{=}66$     & $r^\dagger{=}59$ & $0^{\mathrm{e}}$       \\
TM11 cold key (\classcold{})                       & inv         & $0$                    & $0^{\mathrm{f}}$     & $0$              & $V$                    \\
\bottomrule
\end{tabular}
\par\vspace{3pt}
\footnotesize
\textbf{Zero-loss cells explained:}
$^{\mathrm{a,b,c}}$ no anchor output, hence no descendant-chain
pinning surface; $^{\mathrm{d,e}}$ no partial-unvault-with-revault
primitive (atomic withdrawal only), so no split-exhaustion surface;
$^{\mathrm{f}}$ keyless recovery path, no cold key to compromise.
$^{*}$\catcsfs{}'s TM2 attack requires the cold key; if that key is
compromised the attacker prefers TM11 (full theft). $^{\dagger}$TM3
composes with TM1 to achieve $V$ when fee pinning succeeds.
\end{table}

No row of Table~\ref{tab:pareto-matrix} is dominated in every
column, and TM4's fee-rate-dependence interacts with TM1's
invariance to make aggregate rankings unstable. We return to the
resulting deployment surface in Section~\ref{sec:deployment}.
